\section{Discussion}

% Prøve å overbevise leseren om konklusjonen

% Ta gjerne med referanser

%Nevne at magnetiseringen ideelt sett slutter seg rundt i en sirkel for å minimere magnetostatisk energi ved å unngå frie poler ved kantene
%Hvorfor curl er en god måte å måle hvor perfekte domenene er?

%\subsection{Analysis of results}

The relationship illustrated in Figure \ref{fig:Trend} indicates that as the polygon shape approaches circular symmetry, the magnetic system increasingly favors configurations that minimize stray fields \footnote{Stray magnetic fields are the external magnetic fields generated by internal magnetization of the material. They contribute to magnetostatic energy and can lead to the formation of magnetic domains \cite{definitions2025_strayfield}} by forming energetically favorable closed-loop magnetization patterns.

For all polygons, except the triangles, the average curl is initially high primarily due to strong localized curl along domain walls. %The small overall area amplifies this effect, resulting in a high curl density. 
As the number of sides increases, the curl gradually decreases seemingly linearly, but it fits an exponential decay model as expected, since the curl must eventually stabilize as the number of sides approaches infinity. This trend reflects the increasing geometric resemblance to a circle, where the magnetization becomes more uniformly distributed and the curl is low across the entire structure as seen in Figure \ref{fig:STEM-DPC}.%This trend is further supported by the domain distribution in the figures from heptagon to nonagon, shown in Figure }\ref{fig:STEM-DPC}.

%The discrete nature for the polygons with fewer sides causes the magnetic field to turn drastically along domain walls as can be seen in Figure \ref{fig:triangle}, as opposed to the smoother turns for higher polygon side count, resulting in an overall higher curl per area.

As Figure \ref{fig:Trend} shows, the exponentially decaying trend is consistent and follows the same pattern regardless of the size scale of the studied polygons. The difference is that generally higher curl values were found for smaller polygons. This can be related to a lower number of data points, which increases uncertainty, and that regions where high curl is expected, such as vortices and distinct domain walls, constitute a larger fraction of the total pixel count.

%\subsection{Sources of error}

Several sources of error could have influenced our work with this report. First of all, the resolution of our images is terrible. This is a fault of the TEM laboratory work, and has made it difficult to do sufficient data processing. Second, the FIB milling is not sufficient, as some of the thin film is too thick, and end up influencing and contributing to the domains in some of the structures. We know this because a vortex tendency can be seen for some low-sided polygons if the milled area is seen as a part of the respective polygons. Notice how the magnetic field seems to swirl around a point of interest at the base of the triangle in Figure \ref{fig:triangle}. This may explain why the triangles do not follow the general trend observed, as they, unlike the other polygons, were not fully isolated from the background and are influenced by magnetic interactions.

%A consequence of the poor milling can be seen at the triangles in Figure \ref{fig:STEM-DPC}, for all sizes. The domains are "completed" with the background close to the triangles resulting in a similar domain as those seen for the rectangles, which is also quite visible in Figure \ref{fig:triangle}, where the vector field seems to be rotating around a point outside the structure.

Additionally some are interacting with either the removed film, or potentially the magnetic domains on the film outside the structured region, resulting in a displacement of the domain intersection. If the resulting displacement of the intersection seen in Figure \ref{fig:STEM-DPC} was due to the structures interacting with each other, you would expect the displacement to wane as the structure size decreased, which does not seem to be the case. Another reason for the displacement could be the off angle used for imaging with TEM. Though you would then expect the displacement of the intersection to be consistent between all the shapes, additionally the angle was a consequence of trial and error in the attempt of centering the domain intersection. 